Эффективное функционирование современных научно-информационных систем требует надёжного извлечения и нормализации библиографических ссылок из научных публикаций. В последние годы наблюдается устойчивый рост числа научных публикаций\cite{thelwall2022scopus}, что делает ручную обработку даже одного списка литературы трудоёмкой, а для больших массивов документов -- практически невозможной без автоматизации. Реальная практика работы с библиографией сопряжена со множеством трудностей.

Во первых, для дальнейшего анализа научной публикации необходимо извлечь текст из статьи. В современной практики научные публикации хранятся в формате PDF-файлов, что располагает к совокупности возможных артефактов. PDF -- формат визуальной фиксации, а не семантической разметки, поэтому при извлечении возникают артефакты: разрыв слов, нарушение порядка при двухколонной вёрстке, потеря специальных символов. В случае сканированных документов требуется OCR, который также вносит ошибки, особенно на зашумлённых сканах или с нестандартными шрифтами. Ошибки на этапе извлечения текста носят кумулятивный характер и напрямую влияют на качество последующего выделения библиографии.

Во-вторых, необходимо решить задачу локализации области библиографии в извлеченном текстовом слое документа. Этот этап осложнён рядом деструктивных факторов, препятствующих применению жестких шаблонов поиска. Синтаксический шум: номера страниц, колонтитулы и артефакты PDF-конвертации, -- часто нарушают целостность ключевых слов-маркеров, например, возникновение ложных пробелов вида \guillemotleft Refe~rences\guillemotright{} полностью разрушает регулярные маски. Ситуацию осложняет различное возможное написание заголовков секции, от \guillemotleft Литература\guillemotright, \guillemotleft Bibliography\guillemotright или \guillemotleft Cited works\guillemotright до полного отсуствия их в тексте публикации. При двухколоночном макете элементы верстки могут вклиниваться между основным текстом и финальным списком литературы, а сам раздел библиографии часто предваряет приложения, блоки благодарностей или редакторские примечания. В гуманитарных дисциплинах ссылочная структура и вовсе дробится на постраничные сноски в нижних колонтитулах, что отменяет концепцию сквозной финальной секции.

Третьим шагом конвейера является этап построчной сегментации ссылок из целевой области. Утрата исходного форматирования PDF-документа стирает синтаксические границы разделителей в виде маркированных списков или пустых строк. В результате происходит склеивание служебных маркеров с основным текстом библиографической ссылки. С другой стороны, специфика двухколоночной верстки в условиях некорректного анализа схемы документа провоцирует хаотичное перемешивание текстовых блоков из левой и правой колонок, нарушая линейный порядок следования токенов. Наконец, пространственное смещение границ на этапе локализации приводит либо к необратимой обрезке финальных записей списка литературы, либо к интеграции ложноположительного текстового шума из соседних разделов статьи.

Заключительный этап в конвейере извлечения библиографии -- семантическая сегментация атомарных библиографических областей: автор, заголовок, источник. Задача осложнена кумулятивным эффектом синтаксического \guillemotleft стилевого хаоса\guillemotright{}, внутренней вариативностью сущностей и лавинным накоплением ошибок предыдущих стадий. Экстремальная гетерогенность издательских стандартов кардинально меняет порядок следования полей, типы разделителей и порождает физически вложенные атрибуты, например, неразрывные связки тома, номера и страниц. Высокая энтропия самих текстовых данных -- от наличия в заголовках математических формул, химических нотаций и греческих символов до вариативности антропонимов с национальными частицами (\guillemotleft фон\guillemotright, \guillemotleft де\guillemotright) или инверсией инициалов -- приводит к ложной интерпретации текста как метасимволов-разделителей границ. Наконец, координатные сдвиги, локальные разрывы строк и фантомный шум колонтитулов, унаследованные от первичного PDF-парсинга, полностью разрушают контекстную структуру, лишая стохастические модели возможности корректного распознавания всей последующей цепочки токенов.

Современным деструктивным фактором в задачах анализа научной периодики выступает интеграция генеративного искусственного интеллекта в процесс академического письма, порождающая феномен \guillemotleft генеративных галлюцинаций\guillemotright{} в списках литературы. Широкое вовлечение больших языковых моделей авторами публикаций приводит к автоматическому формированию семантически и синтаксически правдоподобных фантомных библиографических записей, компилирующих реальные антропонимы исследователей и валидные названия авторитетных журналов. Извлечение структурированных полей из таких источников приводит к генерации ложных идентификаторов, таких как: несуществующих DOI, вымышленных томов и страниц, что меняет парадигму обработки данных: традиционный синтаксический анализ становится недостаточным, предопределяя острую необходимость внедрения этапов внешней семантической верификации и перекрестной сверки токенов с глобальными индексами цитирования: Crossref, Scopus, PubMed. Тем не менее, проектирование модулей внешней верификации и интеграция с внешними библиометрическими API выходят за рамки ограничений текущего исследования и рассматриваются как перспективное направление дальнейшей работы.

Несмотря на многолетнюю историю исследований в области работы с библиографией, существующие на данный момент инструменты имеют свои ограничения. Регулярные выражения и другие правило-ориентированные подходы требуют составления шаблонов под каждый новые стиль оформления в журнале. Этот момент лишает универсальности такого подхода. Однако регулярные выражения могут подходить для работы с элементами статьи, похожие во многих изданиях, например название списка литературы. Методы на основе классического машиннного обучения, как CRF или градиентный бустинг, имеют стабильные результаты на обучаемых доменах, но имеют плохую склонность к обобщению при смене стиля оформления. Одним из лучших инструментов для работы с библиографии считается GROBID. Этот инструмент с открытой архитектурой и имеет внутри себя фреймворк для генерации обучающих данных и переобучения моделей.

Для решения этих проблем предложено использовать дообученные БЯМ с использованием LoRA-адаптеров. Актуальность темы исследования обуславливается отсутствием открытого инструмента для разбора библиографии.

Объектом исследования является процесс автоматического извлечения и структурирования библиографической информации из полнотекстовых научных публикаций. Предметом исследования являются методы и алгоритмы для многоуровневого структурирования библиографических записей в условиях стилевого многообразия и шума данных. В частности изучается применимость больших языковых моделей с применением низкоранговой адаптации в области извлечения библиографической информации из научных документов.

Цель работы -- разработка системы для разбора библиографии. Получившийся результат планируется внедрить в ИСАНД\cite{губанов2024информационная}, разрабатываемую в Институте Проблем Управления.

Для достижения поставленной цели необходимо выполнить следующие задачи:
\begin{enumerate}
    \item Провести обзор литературы существующих подходов для разбора библиографии;
    \item Сформулировать формальную постановку задачи извлечения библиографии;
    \item Сформировать набор данных для экспериментов по сравнению различных подходов по извлечению библиографии;
    \item Собрать обучающую выборку для тонкой настройки БЯМ;    
    \item Подобрать большую языковую модель для экспериментов;
    \item Провести сравнительный анализ качества работы различных методов;
    \item Спроектировать гибридную систему для извлечения библиографии;
    \item Интегрировать полученный модуль в систему ИСАНД;
    \item Оценивать вычислительную эффективность разработанного решения.
\end{enumerate}

Для разработки системы потребуется cравнить известные подходы для извлечения библиографии и  выявить наиболее подходящие для каждого подэтапа. В такие подходы входят: регулярные выражения, гистограммные градиентные бустинговые деревья\cite{kopanichuk2025structure}, GROBID\cite{romary2015grobid}, большие языковые модели, дообученные БЯМ и использованием LoRA-адаптеров. В качестве меры качества выступает F1-мера.

Научная значимость и новизна заключается в том, что ранее не проводилось сравнительного анализа извлечения библиографии на русскоязычном наборе данных. Для этого запрошены данные у общероссийского портала Math-Net.Ru. Предоставленный корпус состоял из нескольких тысяч теховских русскоязычных файлов со списками литературы. На полученном наборе данных поставлен эксперимент по выделению с помощью БЯМ библиографических областей.

Практическая значимость исследования заключается в использовании полученного решения в системе ИСАНД. В информационной системе анализа научной деятельности сервис по экстракции библиографических областей может использоваться в двух местах. Во-первых, в личном кабинете на сайте, где пользователь загружает свои публикации. Модуль по извлечению библиографии позволяет тому, кто загрузил свои документы, получить обратный библиографический список. Это перечень заголовков документов, авторов, журналов и других научных агентов, то есть тех, кто создаёт научные публикации. С другой стороны сервис по извлечению библиографии в ИСАНДе может пригодиться для построения кастомизированных библиометрических сетей. Дело в том, что не все публикации проиндексированы в научных базах данных.

Результаты исследования апробированы на конференциях: «Актуальные проблемы семантического анализа данных», «68-я Всероссийская научная конференция МФТИ», «52-е Гагаринские чтения», «Ломоносов-2026».

По итогам работы опубликована статья <<Проблемы извлечения библиографии из научных публикаций>> в журнале «Системы высокой доступности»\cite{system_higher}.

Данная ВКР состоит из введения, заключения и четырех глав. Во второй главе проведен обзор литературы и существующих подходов. В третьей главе сформулирована формальная постановка задачи и описаны метрики качества. В четвертой главе описаны проведенные эксперименты. Пятая глава посвящена обсуждению полученных результатов, анализу архитектуры и интеграции модуля в систему ИСАНД.